\section{问题一模型的建立与求解}

\subsection{模型建立}
将一天按 $10\,$min 离散为 144 个时段，记时间步长 $\Delta t=1/6$\,h。
以各时段的外网购电功率、储能充电功率、储能放电功率及弃光功率为决策变量，
以各时段末的储电量描述储能状态变化。

\textbf{（1）功率平衡约束。} 外网购电、光伏发电及储能放电构成微网的电力供给，
小区负载、储能充电及弃光构成电力消耗。对任一时段 $t$，有
\begin{equation}
p(t)+P_{pv}(t)+u_d(t)=L(t)+u_c(t)+q(t),\qquad t=1,2,\dots,144
\end{equation}
其中 $q(t)$ 为弃光功率，且 $p(t)\ge0$、$q(t)\ge0$；$p(t)\ge0$ 表明微网仅从外网购电，不反向售电。

\textbf{（2）储能运行约束。} 考虑充放电损耗，储能状态转移方程为
\begin{equation}
S(t)=S(t-1)+\eta_c\,u_c(t)\Delta t-\dfrac{u_d(t)\Delta t}{\eta_d},\qquad t=1,2,\dots,144
\end{equation}
其中 $\eta_c=\eta_d=0.9$：充电时仅 $90\%$ 电能存入储能，放电时需多释放 $1/0.9$ 倍电量。
由题目给定的储电量运行范围与充放电功率限额，
\begin{equation}
1200\le S(t)\le 10800,\quad t=0,1,\dots,144;\qquad
0\le u_c(t)\le 5000,\quad 0\le u_d(t)\le 5000 .
\end{equation}
问题一要求 0:00 与 24:00 的储电量相同，设置日循环约束
\begin{equation}
S(144)=S(0).
\end{equation}

\textbf{（3）目标函数。} 第 $t$ 时段购电量为 $p(t)\Delta t$，费用为 $c(t)p(t)\Delta t$，故
\begin{equation}
\min Z=\sum_{t=1}^{144} c(t)\,p(t)\,\Delta t .
\end{equation}

上述模型共 721 个变量（$4\times144$ 个功率变量与 145 个储能状态变量）、289 个等式约束，
目标与约束均为线性，属小型线性规划，可用单纯形法求得全局最优解。

\subsection{模型求解}
将附件 1 的分时电价、小区负载及光伏预测功率代入模型，
采用 Python 调用 HiGHS 线性规划求解器（scipy.optimize.linprog 接口）求解，并用 MATLAB linprog 独立复算，
两者结果逐位一致。购电功率按 $Q_t=p(t)\Delta t$ 换算为各时段购电量输出。

关于日循环边界的取值：附录 1 给定 2025 年 1 月 1 日 0:00 的储电量为 6000\,kWh，
本文在主文中据此取 $S(0)=S(144)=6000$\,kWh；作为对照，将首末储电量改由模型自由确定的情形
（得 $S(0)=S(144)=8550$\,kWh）完整求解并列于附录 A——两情形全天购电量与指定时段购电量完全相同，
费用仅相差 8.35 元（0.024\%），结论对该边界的两种解读均稳健。

\subsection{求解结果与分析}
最优策略下，全天外网购电量为 \textbf{59482.70\,kWh}，全天购电费用为 \textbf{35126.95 元}，
0:00 与 24:00 储电量均为 6000.00\,kWh，全天弃光量为 0。
题目要求的指定时段购电量如表~\ref{tab:q1-plan} 所示，
储能分时段充放电量与首末储电量如表~\ref{tab:q1-soc} 所示，全天调度过程如图~\ref{fig:q1} 所示。

\begin{table}[!htbp]
\caption{微网在指定时间段的购电量及全天的购电量和购电费}\label{tab:q1-plan}
\centering
\begin{tabular}{cccccc}
\toprule
时间段 & 购电量/kWh & 时间段 & 购电量/kWh & 时间段 & 购电量/kWh \\
\midrule
10:00--10:10 & 0.0000 & 12:00--12:10 & 486.4029 & 14:00--14:10 & 0.0000 \\
16:00--16:10 & 394.9315 & 18:00--18:10 & 636.9826 & 20:00--20:10 & 0.0000 \\
\multicolumn{3}{l}{全天购电量/kWh：59482.70} & \multicolumn{3}{l}{全天购电费/元：35126.95} \\
\bottomrule
\end{tabular}
\end{table}

\begin{table}[!htbp]
\caption{储能设备在指定时间段的充放电量及 0:00 和 24:00 的储电量}\label{tab:q1-soc}
\centering
\begin{tabular}{cccccc}
\toprule
时间段 & 充电量/kWh & 放电量/kWh & 时间段 & 充电量/kWh & 放电量/kWh \\
\midrule
0:00--4:00   & 4500.0000 & 0.0000    & 4:00--8:00   & 833.3333  & 5947.4197 \\
8:00--12:00  & 3954.6310 & 2121.4184 & 12:00--16:00 & 6119.3685 & 91.1014 \\
16:00--20:00 & 0.0000    & 5068.6869 & 20:00--24:00 & 5333.3333 & 3571.3131 \\
\multicolumn{3}{l}{0:00 储电量/kWh：6000.0000} & \multicolumn{3}{l}{24:00 储电量/kWh：6000.0000} \\
\bottomrule
\end{tabular}
\end{table}

\begin{figure}[!htbp]
\centering
\includegraphics[width=0.96\textwidth]{Q1_CaseB.pdf}
\caption{问题一最优日前调度结果（$S(0)=S(24{:}00)=6000$\,kWh）：(a) 分时电价；(b) 购电功率与小区负载；
(c) 储能充放电功率（向上为充电、向下为放电）；(d) 储电量轨迹（虚线为运行上下限 1200/10800\,kWh）}
\label{fig:q1}
\end{figure}

由图~\ref{fig:q1} 可见，储能充放电行为与分时电价、光伏出力高度协同：
（i）夜间低谷（22:00 至次日 6:00，电价约 0.37$\sim$0.44 元/kWh）低价购电满足负载并对储能充电，
0:00--4:00 与 20:00--24:00 两段合计充电 9833.33\,kWh；
（ii）早高峰（6:00--10:00，电价 0.69$\sim$1.17 元）储能大功率放电（4:00--8:00 段放电 5947.42\,kWh），压低高价购电；
（iii）中午光伏富余且电价处于日内低谷，富余光伏全部存入储能，并以低价购电充满充电通道（12:00--12:10 购电 486.40\,kWh），
供晚高峰使用——存储套利的经济性条件 $0.9\times0.9\times1.35>0.44$ 成立；
（iv）傍晚峰段（18:00--21:00，电价 1.23$\sim$1.40 元）储能能量稀缺，
优先留给最高价的 18:50--20:50 放电（该时段购电为 0），其余峰段由购电覆盖净负载。
全天共有约 20.7\,MWh 电量经储能周转，每 1\,kWh 夜间电量经储能（往返效率 0.81）在峰时释放净收益约
$0.81\times1.35-0.43\approx0.66$ 元；与不配置储能的情形（48052.05 元，见第九章）相比，
储能使全天购电费下降 \textbf{12925.10 元（26.9\%）}。
